\chapter{Full Adders}
\label{chap:full-adders}

\chapterintro{A full adder extends a half adder by accepting a carry-in bit. This allows many adders to be chained together to add multi-bit binary numbers.}

\learningobjectives{
\begin{itemize}
    \item Explain carry-in and carry-out.
    \item Build a full adder from half adders.
    \item Add multi-bit binary numbers conceptually.
    \item Implement a ripple-carry adder in Python.
\end{itemize}}

\section{Why Half Adders Are Not Enough}

A half adder can add two bits, but multi-bit addition also needs to include a carry from the previous column.

A full adder accepts three inputs:

\begin{itemize}
    \item \(A\)
    \item \(B\)
    \item \(C_{in}\), the carry-in
\end{itemize}

It produces:

\begin{itemize}
    \item \(S\), the sum bit
    \item \(C_{out}\), the carry-out bit
\end{itemize}

\section{Full Adder Equations}

The sum bit is:

\[
S=A\operatorname{XOR}B\operatorname{XOR}C_{in}
\]

The carry-out is true when at least two inputs are true:

\[
C_{out}=(A\operatorname{AND}B)\operatorname{OR}(C_{in}\operatorname{AND}(A\operatorname{XOR}B))
\]

\section{Python Implementation}

\begin{lstlisting}
def full_adder(a: int, b: int, carry_in: int) -> tuple[int, int]:
    if a not in (0, 1) or b not in (0, 1) or carry_in not in (0, 1):
        raise ValueError("inputs must be 0 or 1")
    total = a + b + carry_in
    return total % 2, total // 2
\end{lstlisting}

\section{Ripple-Carry Addition}

A ripple-carry adder chains full adders from the least significant bit to the most significant bit. Each carry-out becomes the next carry-in.

\section{Exercises}

\begin{exercise}
For \(A=1\), \(B=1\), and \(C_{in}=1\), compute \(S\) and \(C_{out}\).
\end{exercise}

\begin{solution}
The total is \(1+1+1=3\), which is \(11_2\). Therefore \(S=1\) and \(C_{out}=1\).
\end{solution}

\section{Portfolio Milestone}

Implement a ripple-carry adder that accepts two binary strings and returns their binary sum.

\section{Summary}

Full adders make multi-bit binary addition possible. They are a fundamental building block of arithmetic logic units.
